\documentclass[10pt,a4paper]{amsart}
\usepackage[margin=19mm]{geometry}
\usepackage{amsmath,amssymb,mathtools}
\usepackage[scr=rsfs,cal=euler]{mathalfa}
\usepackage{xfrac}
\usepackage[unicode,hidelinks,bookmarks=false]{hyperref}
\newcommand{\ie}{i.e.\ }
\newcommand{\cf}{cf.\ }
\newcommand{\resp}{respectively\ }
\newcommand{\Subsection}{\subsection}
\providecommand{\nobreakdash}{\nobreak\hbox{-}}
\setlength{\emergencystretch}{2em}
\multlinegap0pt
\usepackage{xcolor}
\usepackage{amssymb}
\usepackage{tikz}
\usepackage[normalem]{ulem}
\newtheorem{cor}{Corollary}
\newtheorem{lem}{Lemma}
\newcommand{\N}{\mathbb{N}}
\newcommand{\bfs}{{\boldsymbol{\sl{s}}}}
\newcommand{\calA}{\mathcal{A}}
\newcommand{\ol}{\overline}
\title{Finite and Symmetric Multiple $T$-Values}
\author{Aaron Cheng, Jianqiang Zhao}
\date{Source: arXiv:2608.27541v1 (CC BY 4.0)}
\begin{document}
\maketitle

\noindent\emph{Source and licence.} Finite and Symmetric Multiple $T$-Values. Version arXiv:2608.27541v1 (CC BY 4.0), distributed by arXiv under the Creative Commons Attribution 4.0 International licence, http://creativecommons.org/licenses/by/4.0/, as recorded in the arXiv licence metadata for this exact version and shown on the abstract page https://arxiv.org/abs/2608.27541v1. Full licence notice: \texttt{licenses/arxiv-2608.27541-CC-BY-4.0.txt}.\par\smallskip

\noindent\textit{Editorial scope.} Counted unit: the article's lem lem:S2(x) (source0.tex lines 936--943). The article's complete original proof of it is delivered with it (source0.tex lines 945--975, one \texttt{proof} environment of 31 lines). No mathematics is written, repaired or re-derived: the statement and the proof are the article's own lines, in the article's own notation, and no retained line is substituted or re-flowed.  Retained uncounted as the source-local context the proof needs: lines 498--510.  Nothing was omitted from any retained span, so the package carries no omission record.  No retained line contains a citation, so the wrapper carries no bibliography and no bibliography entry.  No retained line contains a figure environment, an image-inclusion command or drawing code, so no image is delivered and the package declares no asset dependency.  Editorial formatting only, in addition: an AMS \texttt{amsart} preamble that restates the article's own declarations and macros used by the retained lines, namely the environments \texttt{cor}, \texttt{lem} on separate counters, together with the standard packages those lines need.  The retained environments print in this document as Corollary 1, Lemma 1 in order of appearance, whereas the article prints the same items with the numbering of its own full document.  The complete native source of the article as distributed in the arXiv e-print for this exact version is bundled unchanged as \texttt{sources/208720/source0.tex}.
\begin{cor}\label{cor:unitT-oddWt}
Let $k$ be a positive integer. Then
\begin{equation*}
T_\calA(\{1\}^k)=T_\calA(k)=-\zeta_\calA(\ol{k})=
\left\{
  \begin{array}{ll}
    0, & \hbox{if $k$ is even;} \\
    2\texttt{q}_2, & \hbox{if $k=1$;} \\
    2(1-2^{1-k})\beta_k, & \hbox{if $k$ is odd and $k\ge 3$.}
  \end{array}
\right.
\end{equation*}
\end{cor}

\begin{lem}\label{lem:S2(x)}
For any index $\bfs=(s_1,\dots,s_d)\in\N^d$ and $k\in\N$, define $\bfs+k=(s_1,\dots,s_{d-1},s_d+k).$
Then, for any prime $p>k+1$, we have
\begin{equation*}
S_\calA(\bfs,k+1)_p\equiv -T_\calA(k+1)_p \cdot T_\calA(\bfs)_p
+\sum_{j=0}^{p-k-2}\binom{p-k-1}{j} 2^jB_j\Big(\frac12\Big)\frac{T_{\calA}(\bfs+k+j)_p }{p-k-1}\pmod{p}.
\end{equation*}
\end{lem}

\begin{proof}
Set $N=p-k-1\ge 1$ for simplicity. Then
\begin{align*}
\ S_\calA(\mathbf s,k+1)_p&
\equiv\sum_{\substack{p>r_1>\cdots >r_{d+1}>0 \\ r_{j}\equiv d+1-j\pmod 2}} \frac{2^{d+1}}{r_1^{s_1}\cdots r_{d}^{s_d}r_{d+1}^{k+1}} \\
&\equiv\sum_{\substack{p>r_1>\cdots >r_{d}>0\\r_j\equiv d+1-j\pmod 2}} \frac{2^{d+1}}{r_1^{s_1}\cdots r_{d}^{s_d}}\sum_{1\le r_{d+1}'\le (r_{d}-1)/2}\frac{1}{(2r_{d+1}')^{k+1}} \qquad(\text{since $r_{d}$ is odd}) \\
&\equiv\sum_{\substack{p>r_1>\cdots >r_{d}>0\\r_j\equiv d+1-j\pmod 2}} \frac{2^{d+1}}{r_1^{s_1}\cdots r_{d}^{s_d}} \cdot\frac{2^{-k-1}}{N}\sum_{j=0}^{N-1}(-1)^j\binom{N}{j}B_j\left(\frac{r_{d}-1}{2}\right)^{N-j}.
\end{align*}
Expanding the inner sum by adding and subtracting the extra term corresponding to $j=N$, we have
\begin{align*}
&\,\sum_{j=0}^{N}(-1)^j\binom{N}{j}B_j\left(\frac{r_{d}-1}{2}\right)^{N-j}-(-1)^N B_N\\
=&\,
\sum_{j=0}^{N}(-1)^j\binom{N}{j}B_j2^{j-N}\sum_{i=0}^{N-j}\binom{N-j}{i}(-1)^{N-j-i}r_{d}^{i}-(-1)^N B_N\\
=&\, \sum_{0\le i+j\le N} (-1)^{k+i}\binom{N}{j}\binom{N-j}{i}B_j2^{j-N}r_{d}^{i} -(-1)^N B_N\\
=&\, \sum_{0\le i+j\le N} (-1)^{k+i}\binom{N}{i}\binom{N-i}{j}B_j2^{j-N}r_{d}^{i} -(-1)^N B_N\\
=&\, \sum_{j=0}^{N}2^{-j}\Bigg\{\sum_{i=0}^{N-j}\binom{N-j}{i}B_i2^{i-(N-j)}\Bigg\}(-1)^{k+j}\binom{N}{j}r_{d}^{j}-(-1)^N B_N \quad(\text{swapping $i$ and $j$})\\
=&\, \sum_{j=0}^{N}2^{-j}B_{N-j}\Big(\frac12\Big)(-1)^{k+j}\binom{N}{j}r_{d}^{j}-(-1)^N B_N\\
=&\, \sum_{j=0}^{N}2^{-(N-j)}B_{j}\Big(\frac12\Big)(-1)^{k+N-j}\binom{N}{j}r_{d}^{N-j}-(-1)^N B_N \\
=&\, \sum_{j=0}^{N-1}2^{-(N-j)}B_{j}\Big(\frac12\Big)(-1)^{k+N-j}\binom{N}{j}r_{d}^{N-j}
+(-1)^k\Big(B_{N}\Big(\frac12\Big)- B_N\Big) .
\end{align*}
From this we have
\begin{align*}
S_\calA(\bfs,k+1)_p\equiv&\,
\frac{2^{-k}}{p-k-1}(-1)^k(2^{1+k}-2)B_{p-k-1}\sum_{\substack{p>r_1>\cdots >r_{d}>0\\r_j\equiv d+1-j\pmod 2}}\frac{2^{d}}{r_1^{s_1}\cdots r_d^{s_{d}}} \\ &\,+\frac{2^{-k}}{p-k-1}\sum_{j=0}^{p-k-2}2^{k+j}B_j\left(\frac12\right)\binom{p-k-1}{j}
\sum_{\substack{p>r_1>\cdots >r_{d}>0\\r_j\equiv d+1-j\pmod 2}}\frac{2^{d}}{r_1^{s_1}\cdots r_{d-1}^{s_{d-1}}}r_d^{p-k-1-j-s_{d}}\\
&\equiv (-1)^{k-1}T_\calA(k+1)_pT_\calA(\bfs)_p
+\frac{1}{p-k-1}\sum_{j=0}^{p-k-2}2^jB_j\left(\frac12\right)\binom{p-k-1}{j} T_{\calA}(\mathbf s+k+j)_p,
\end{align*}
by Corollary~\ref{cor:unitT-oddWt}. We can even replace $(-1)^{k-1}$ by $-1$ since this term is 0 when $k$ is odd.
\end{proof}

\end{document}
